%%
%% template.tex
%%
%% Starting point for a new DUNE Interface Control Document.  Copy this
%% file (and DUNEICD.cls, dune-icd-changes.sty, resolve_changes.py, Logos/,
%% .gitignore from this _common/ directory) into a new project directory,
%% rename it to dune_icd_daq_<interface>.tex, and fill in the placeholders
%% below.
%%
%% Build with:  latexmk -pdf dune_icd_daq_<interface>.tex
%%
\documentclass[letterpaper]{DUNEICD}

%--------------------------------------------------------------------------
% Editorial markup (dune-icd-changes package)
%--------------------------------------------------------------------------
%
%   \added[id=WK]{new text}
%   \deleted[id=WK]{old text}
%   \replaced[id=WK]{new text}{old text}
%   \comment[id=WK]{question or note for reviewers}
%   \begin{addedblock}[id=WK] ... \end{addedblock}
%   \begin{deletedblock}[id=WK] ... \end{deletedblock}
%
% See dune-icd-changes.sty for the full command reference and the
% rationale for addedblock/deletedblock (structural changes that \added/
% \deleted cannot express, since they are LR-mode commands).
%
% >>> BEFORE SUBMITTING TO EDMS: swap the two \usepackage lines below and
%     comment out \listofchanges at the end of this document. <<<
%
\usepackage{dune-icd-changes}          % --- DRAFT / REVIEW mode (markup visible) ---
%\usepackage[final]{dune-icd-changes}  % --- SUBMISSION mode (markup silently resolved) ---

% Reviewers / editors.  Colors come from xcolor's dvipsnames set, which the
% DUNEICD class already loads.  One \definechangesauthor per person who
% might leave markup; the [id] used in \added/\comment/etc must match.
\definechangesauthor[name={Your Name},color=blue]{YN}

%--------------------------------------------------------------------------
% Document metadata
%--------------------------------------------------------------------------
\DocTitle{Full Document Title}
\ShortDocTitle{DUNE ICD: Short Title}
\DocDescription{One or two sentence description of what this document
  covers, for the cover page.}
\EDMSDocId{XXXXXXX}
\EDMSDocUrl{https://edms.cern.ch/document/XXXXXXX}
\DocVersion{0.1}
\CreationDate{\today}
\ModifiedDate{\today}
\DocAuthors{Author One (Inst)\\ Author Two (Inst)}
\DocCheckedBy{Reviewer One (Inst)\\ Reviewer Two (Inst)}
\DocApprovedBy{Approver One (Inst)\\ Approver Two (Inst)}
\DistributionList{\parbox[t]{0.7\textwidth}{\raggedright
  \href{mailto:list@example.org}{list@example.org}}}

\begin{document}

\makeICDcover

\begin{revisionhistory}
  \revEntry{0.1}{\today}{Author}{Initial draft}
\end{revisionhistory}

\tableofcontents
\clearpage

%--------------------------------------------------------------------------
\section{Introduction}
\label{sec:introduction}
%--------------------------------------------------------------------------

Describe the purpose and scope of this interface document: which two
consortia/groups it connects, and what it is meant to formalize.

%--------------------------------------------------------------------------
\section{Section Title}
\label{sec:section-one}
%--------------------------------------------------------------------------

Body text.  Use \texttt{item\_dense}/\texttt{enum\_dense} for compact lists:

\begin{item_dense}
  \item First point.
  \item Second point.
\end{item_dense}

An item/responsibility/status table, e.g. for open action items or
deliverables (see DUNEICD.cls for the \texttt{icdtable}/\texttt{icdrow}
definitions):

\begin{icdtable}{Activity}
  \icdrow{1}{Description of the first item.}{Responsible party}{\textbf{OPEN}}
  \icdrow{2}{Description of the second item.}{Responsible party}{Planned}
\end{icdtable}

%--------------------------------------------------------------------------
\section{References}
\label{sec:references}
%--------------------------------------------------------------------------

List related EDMS documents and other references here.

%--------------------------------------------------------------------------
% Only include \appendix and the sections below if this document needs
% appendices.  Use plain \appendix + \section{...} -- NOT a hand-rolled
% "APPENDIX A" section -- so lettering, ToC entries, and \ref/\label all
% work correctly. Referenced elsewhere as e.g. Appendix~\ref{app:example}.
%--------------------------------------------------------------------------
% \clearpage
% \appendix
% \section{Appendix Title}
% \label{app:example}
%
% Appendix body.

%--------------------------------------------------------------------------
% Summary of editorial markup -- comment out for the EDMS submission build.
%--------------------------------------------------------------------------
\clearpage
\listofchanges

\end{document}
